% Auxiliary finite-leg form of the charge detection calculation in SCP10,
% local source lines 2464--2486. Checked auxiliary statements.
\begin{definition}[Shared-leg coordinates and literal charged columns]
    \label{def:peps_regular_two_site_charge_coordinates}
    Let $I,J$ be finite incident-leg sets, with distinguished shared legs $i_0,j_0$.
    Write $I_0=I\setminus\{i_0\}$ and $J_0=J\setminus\{j_0\}$. Splitting the
    shared labels gives the bijection
    \[
        G^I\times G^J\simeq(G\times G)\times(G^{I_0}\times G^{J_0}).
    \]
    If $\iota_I(k,\theta_I)$ inserts $k$ at the shared leg, the actual canonical
    charged column is
    \[
        C_{\chi,p,\theta}(\alpha,\beta)=\sum_{k\in G}\chi(pk)
        P_I(\alpha,\iota_I(k,\theta_I))P_J(\beta,\iota_J(k,\theta_J)).
    \]
    Here $P_I,P_J$ are the local regular averaging projectors.
    \lean{TNLean.PEPS.regularSharedLegEquiv,TNLean.PEPS.regularTwoSiteChargeColumn}
    \uses{def:peps_regular_site_gram}
    \leanok
\end{definition}

\begin{theorem}[Actual two-site charge expansion]
    \label{thm:peps_regular_two_site_charge_expansion}
    Put $N=|G|$ and
    $A_{\chi,p;x,y}(r,s)=\mathbf1_{s=yx^{-1}r}\chi(px^{-1}r)$.
    The literal contraction satisfies
    \[
        C_{\chi,p,\theta}(\alpha,\beta)=N^{-2}\sum_{x,y\in G}
        A_{\chi,p;x,y}(\alpha_{i_0},\beta_{j_0})
        \mathbf1_{\alpha|_{I_0}=x\theta_I}
        \mathbf1_{\beta|_{J_0}=y\theta_J}.
    \]
    It is also $(P_I\otimes P_J)$ applied to the character-weighted shared bond,
    with every other boundary label retained.
    \lean{TNLean.PEPS.regularTwoSiteChargeColumn_apply,TNLean.PEPS.regularTwoSiteChargeColumn_eq_projector}
    \uses{def:peps_regular_two_site_charge_coordinates,def:peps_regular_charge_coordinates}
    \leanok
\end{theorem}
\begin{proof}\leanok
    Expand both local averages and sum the shared label. For each $x,y$ its value
    is forced to be $x^{-1}\alpha_{i_0}$, giving the displayed charge matrix.
\end{proof}

\begin{definition}[Shared-leg charge detector]
    \label{def:peps_regular_two_site_charge_detector}
    Let $D_\chi$ be the orthogonal projector onto the span of all translated charge
    matrices $A_{\chi,p;x,y}$. In the shared-leg coordinates define
    $\widehat D_\chi=D_\chi\otimes I$ on the remaining labels, and transport it
    back by the preceding bijection.
    \lean{TNLean.PEPS.regularTwoSiteChargeDetector}
    \uses{def:peps_regular_two_site_charge_coordinates}
    \leanok
\end{definition}

\begin{theorem}[Canonical charge projection and local averages]
    \label{thm:peps_regular_two_site_charge_detector}
    The operator $\widehat D_\chi$ is a positive Hermitian projection, commutes with
    $P_I\otimes P_J$, and satisfies
    \[
        \widehat D_\chi C_{\chi,p,\theta}=C_{\chi,p,\theta}.
    \]
    If $\chi,\psi$ are distinct irreducible characters and $\chi$ is the character
    of a finite-dimensional unitary representation, then
    $\widehat D_\chi C_{\psi,p,\theta}=0$ for every $p,\theta$.
    Its kernel is invariant under independent translations of the two vertices.
    \lean{TNLean.PEPS.regularTwoSiteChargeDetector_properties,TNLean.PEPS.regularTwoSiteChargeDetector_column,TNLean.PEPS.regularTwoSiteChargeDetector_other_column,TNLean.PEPS.regularTwoSiteChargeDetector_translate_entries,TNLean.PEPS.regularTwoSiteChargeDetector_commute_projector}
    \uses{thm:peps_regular_two_site_charge_expansion,def:peps_regular_two_site_charge_detector}
    \leanok
\end{theorem}
\begin{proof}\leanok
    The shared-pair projector commutes with independent endpoint translations.
    Tensoring with the identity preserves this commutation and the projection
    properties. Averaging gives commutation with the local projectors, and the
    literal expansion gives the two eigenvalue formulas.
\end{proof}

\begin{definition}[Original physical charged column]
    \label{def:peps_regular_two_site_physical_charge_column}
    For local tensors $a:G^I\to\mathcal H_A$ and $b:G^J\to\mathcal H_B$, define
    \[
        F_{\chi,p,\theta}(s,t)=\sum_{k\in G}\chi(pk)
        a(\iota_I(k,\theta_I),s)b(\iota_J(k,\theta_J),t).
    \]
    This is the actual contraction of the two original physical tensors along the
    character-weighted shared bond.
    \lean{TNLean.PEPS.regularTwoSitePhysicalChargeColumn}
    \uses{def:peps_regular_two_site_charge_coordinates}
    \leanok
\end{definition}

\begin{theorem}[Two-site physical charge measurement]
    \label{thm:peps_regular_two_site_physical_charge_measurement}
    Assume the two original site maps are regular $G$-isometric, with their positive
    local isometry factors. Their product $T$ satisfies
    $F_{\chi,p,\theta}=T C_{\chi,p,\theta}$. There is a positive Hermitian projection
    $Q_\chi$ on $\mathcal H_A\otimes\mathcal H_B$, chosen before all $p,\theta$, such that
    \[
        Q_\chi T=T\widehat D_\chi,
        \qquad Q_\chi F_{\chi,p,\theta}=F_{\chi,p,\theta}.
    \]
    For a selected irreducible unitary character $\chi$, the pair
    $(Q_\chi,I-Q_\chi)$ is a complete projective measurement. It has eigenvalue
    zero on every $F_{\psi,p,\theta}$ for an inequivalent irreducible character
    $\psi$. Positivity, completeness, and the physical coefficient identities are
    derived from the original local isometry assumptions.
    \lean{TNLean.PEPS.exists_physicalProjectionFamily,TNLean.PEPS.regularTwoSitePhysicalChargeColumn_eq_image,TNLean.PEPS.exists_regularTwoSitePhysicalChargeProjector,TNLean.PEPS.exists_regularTwoSitePhysicalChargeMeasurement}
    \uses{thm:peps_regular_two_site_charge_detector,def:peps_regular_two_site_physical_charge_column}
    \leanok
\end{theorem}
\begin{proof}\leanok
    Derive $T^\dagger T=c(P_I\otimes P_J)$ with $c>0$ from the two local Gram
    identities. Then $Q_\chi=c^{-1}T\widehat D_\chi T^\dagger$ is the required
    projection. Its complement supplies the second outcome, including unused
    physical directions.
\end{proof}

These are auxiliary finite-leg identities for the charge detection argument in
SCP10, lines 2464--2486. No nonvanishing assertion is made when a site has no
remaining boundary leg. A source charge detector on a prescribed lattice block
still requires its actual geometry and nondegenerate contraction comparison;
no six-spin creation or parent-Hamiltonian assertion is included here.


\begin{theorem}[Nonzero original two-site charge columns]
    \label{thm:peps_regular_two_site_charge_nonzero}
    Suppose $I_0,J_0$ are nonempty. At the reconstructed diagonal shared row $k$,
    \[
        C_{\chi,p,\theta}(\iota_I(k,\theta_I),\iota_J(k,\theta_J))
        =|G|^{-2}\chi(pk).
    \]
    Thus $\chi(1)\ne0$ implies that every canonical charged column is nonzero.
    For regular $G$-injective original physical site maps, the actual physical
    charged columns are also nonzero for every internal and boundary label.
    For every actual regular irreducible charge label, the condition
    $\chi(1)\ne0$ follows from the positive dimension of its irreducible
    representative and need not be supplied separately.
    \lean{TNLean.PEPS.regularTwoSiteChargeColumn_diagonal,TNLean.PEPS.regularTwoSiteChargeColumn_ne_zero,TNLean.PEPS.regularTwoSitePhysicalChargeColumn_ne_zero,
        TNLean.PEPS.regularTwoSitePhysicalChargeColumn_ne_zero_of_mem_regularChargeLabels}
    \uses{thm:peps_regular_two_site_charge_expansion,def:peps_regular_two_site_physical_charge_column,thm:peps_regular_charge_labels}
    \leanok
\end{theorem}
\begin{proof}
    \leanok
    The remaining legs force both translations to be the identity. Choose
    $k=p^{-1}$, then apply the derived local $G$-injective inverses to the actual
    physical contraction.
\end{proof}
